\begingroup
\footnotesize
\begin{longtable}{L{2.6cm}L{3.9cm}L{4.4cm}L{\dimexpr\textwidth-10.9cm-8\tabcolsep\relax}}
\caption{Cambios entre el diseño de referencia (primer corte) y el prototipo implementado.}\label{tab:cambios_prototipo}\\
\toprule
\textbf{Elemento} & \textbf{Diseño de referencia (matrices Pugh)} & \textbf{Prototipo implementado} & \textbf{Motivo}\\
\midrule
\endfirsthead
\multicolumn{4}{l}{\small \tablename\ \thetable{} -- continuación}\\
\toprule
\textbf{Elemento} & \textbf{Diseño de referencia (matrices Pugh)} & \textbf{Prototipo implementado} & \textbf{Motivo}\\
\midrule
\endhead
\midrule
\multicolumn{4}{r}{\small Continúa en la siguiente página}\\
\endfoot
\bottomrule
\endlastfoot
Esquema hidráulico & Lazo cerrado con recirculación (A2) & Paso único hacia el tanque de efluentes (TK-103), vaciado manualmente al drenaje entre fases & Simplifica la arquitectura hidráulica y evita la acumulación de residuos en el circuito de preparación (Sección~\ref{sec:caracterizacion})\\
Temperatura de consigna & 50\,°C (preenjuague) y 40\,°C (lavado) & 35\,°C en ambas fases & Margen frente al límite de 40\,°C del sello mecánico de la bomba\\
Bomba de impulsión & Truper BOAC-1/2 EXPERT, centrífuga & Bomba periférica (P-102) & Costo y disponibilidad\\
Agitación & Motorreductor de 100\,rpm con agitador & Bomba sumergible de 22\,W (P-103), activa antes del calentamiento & Evita sumergir un eje en el tanque y sellar su paso por la tapa, que añadía un punto de fuga y piezas móviles en contacto con el fluido\\
Protección térmica de la resistencia & Termostato bimetálico en serie & Sin termostato; corte del calentador por programa del PLC (umbral de 35\,°C) & Descartado en este corte por costo y por simplificación del diseño\\
Conmutación de potencia & Un SSR de 40\,A para la resistencia & SSR de 60\,A para la resistencia y SSR de 40\,A para la bomba de impulsión & El SSR de 60\,A da margen frente a los 16\,A de la resistencia; el de 40\,A estaba disponible en la universidad\\
Disipación de los SSR & Disipador $\leq 1{,}2$\,°C/W y ventilación forzada & Solo la placa metálica propia de cada SSR & Costo; la temperatura de los SSR se vigila durante el calentamiento\\
Accionamiento de la bomba & Contactor y guardamotor & SSR de 40\,A e interruptor automático & Simplificación y costo: la bomba se conmuta con el SSR de 40\,A\\
Electroválvulas & Una electroválvula de llenado & Dos electroválvulas: llenado (XV-102), instalada, y descarga (XV-103), de solenoide y de acción directa, pendiente de compra & Automatizar desde el PLC la descarga de TK-102 hacia la bomba de impulsión\\
Válvulas de bola (regulación y desvío) y válvula antirretorno & Incluidas en el diseño & No se instalan & Descartadas en este corte por costo y por simplificación del diseño hidráulico\\
Filtro Y & Incluido en el diseño & No se usa & Costo y simplificación del circuito de paso único\\
Manómetro & Incluido en el diseño & Aún no instalado & Pendiente de compra\\
Bomba dosificadora & Peristáltica de 24\,V con tubo Norprene & Peristáltica de 12\,V (P-101) alimentada con un convertidor LM2596 & Disponibilidad y costo\\
Sensor de nivel & Dos flotadores reed de 115\,mm & Un flotador doble (LSL-101 y LSH-101) & Mide el nivel bajo y el alto en una sola unidad y con una sola perforación\\
Tanque de mezcla & Acero inoxidable con buje soldado para la resistencia & Tanque de plástico de 50\,L; la resistencia atraviesa la pared con contratuerca y empaques & Tras un percance con el tanque metálico durante el montaje, el tanque de plástico resultó más fácil de perforar y sellar; además, es de menor costo\\
Tubería & Acero inoxidable / CPVC & PVC de 1 y 3/4 de pulgada & Costo y facilidad de montaje, dentro de la temperatura de operación (máximo 35\,°C)\\
Supervisión remota & MQTT con ThingsBoard autohospedado & Raspberry Pi con Node-RED (protocolo S7), API en AWS e interfaz en Streamlit & Arquitectura en definición; propuesta en el Anexo~\ref{anx:arquitectura_com}, a confirmar en el tercer corte\\
\end{longtable}
\endgroup
